\documentclass[11pt,a4paper]{article}
\usepackage[T1]{fontenc}
\usepackage[utf8]{inputenc}
\usepackage{lmodern,microtype}
\usepackage[margin=24mm,headheight=14pt]{geometry}
\usepackage{amsmath,amssymb,amsthm,mathtools}
\usepackage{booktabs,longtable,array,enumitem,xcolor}
\usepackage{needspace}
\usepackage{hyperref}
\hypersetup{colorlinks=true,linkcolor=blue!45!black,urlcolor=blue!50!black,citecolor=blue!40!black,
 pdftitle={Smaller Diophantine forms of Con(ZFC)},pdfauthor={Research collaboration for ProveIt}}
\usepackage{fancyhdr}
\pagestyle{fancy}\fancyhf{}\fancyhead[L]{\small Diophantine forms of ZFC consistency}\fancyhead[R]{\small 22 October 2026}\fancyfoot[C]{\thepage}
\setlength{\emergencystretch}{2em}
\setlist{itemsep=3pt,topsep=5pt}
\newtheorem{theorem}{Theorem}[section]
\newtheorem{lemma}[theorem]{Lemma}
\newtheorem{proposition}[theorem]{Proposition}
\theoremstyle{remark}\newtheorem{remark}[theorem]{Remark}
\newcommand{\Pos}{\mathbb Z_{>0}}
\newcommand{\NN}{\mathbb N}
\newcommand{\ZZ}{\mathbb Z}
\newcommand{\QQ}{\mathbb Q}
\newcommand{\Con}{\operatorname{Con}}
\newcommand{\pc}{\operatorname{pc}}
\newcommand{\ord}{\operatorname{ord}}
\newcommand{\src}[1]{{\small\path{#1}}}
\newcommand{\repo}{https://github.com/VladimirReshetnikov/ProveIt}
\newcommand{\pin}{3b201cdea32d2c816f5f3d20734051f25fd3ef14}
\title{\textbf{Smaller Diophantine forms of \(\Con(\mathrm{ZFC})\)}\\[5pt]
\large Lower degree, fewer unknowns, and smaller fixed coefficients}
\author{Research synthesis prepared for ProveIt}
\date{22 October 2026}
\begin{document}
\maketitle
\begin{abstract}
We combine the eight fixed-input packages in ProveIt's
\texttt{con-zfc-diophantine-forms} collection with its ordinary strong-Pell,
signed-quotient, and quadratic-root-elimination arguments. The resulting
improvements are explicit arithmetic circuits: a four-equation system in
15 positive unknowns, of maximum degree 28 and 63 operations; its single
degree-56 sum of squares, at 74 operations; single polynomials in 14 positive
unknowns at 73 operations/degree 342 and 75 operations/degree 293; a quartic
in 33 positive unknowns; and equations in 33 unrestricted integer unknowns.
A separate finite-window reduction replaces nine-cell windows by projected
crosses, reducing the exact coefficient-width parameter from \(5^{67}\) to
\(5^{28}\). Each completed construction has a solution exactly when the
literal K215 machine halts on blank tape. The established syntactic
calibration therefore gives, by mathematics formalizable in ZFC, that
unsolvability is equivalent to \(\Con(\mathrm{ZFC})\). This packet contains
proofs, complete circuits, fixed coefficient definitions, and reproducible
finite checks. It is not a new end-to-end Lean or Rocq verification, and it
claims no global minimum.
\end{abstract}
\tableofcontents
\newpage

\section{Results, comparison criteria, and scope}

\subsection{The theorem being constructed}

Fix a conventional effective proof predicate for ZFC, with consistency
meaning that no finite proof of falsity exists. All four coefficient ports
in this report denote \emph{fixed integers}, supplied by the terminating
recipe in Section~\ref{sec:coefficients}. There is no quantified program
index, input, exponent, or computation predicate in the final equations.

\begin{theorem}[Fixed Diophantine consistency forms]\label{thm:main}
For each polynomial or system in Table~\ref{tab:new}, with its declared
domain and the fixed coefficients of Section~\ref{sec:coefficients}, ZFC
proves
\[
 \text{the equation or system has no solution}
 \quad\Longleftrightarrow\quad \Con(\mathrm{ZFC}).
\]
The same arithmetic reductions also apply with the original valid
inputless coefficients of package 06. The theorem does not assume
\(\Con(\mathrm{ZFC})\).
\end{theorem}

The proof has three components. The new algebraic reductions preserve
existence of a solution. The new coefficient recipe encodes exactly the
same blank K215 computation problem. The repository's K215 calibration
identifies blank halting with a ZFC refutation. The last step is used as
ordinary formalizable mathematics, not as an appeal to the soundness of
Lean from inside ZFC; see Section~\ref{sec:logic}.

\begin{table}[ht]
\centering\small
\begin{tabular}{@{}p{55mm}rrrr@{}}
\toprule
Construction & Unknowns & Equations & Degree & Ops\\
\midrule
Inherited \(P_{67}\), new coefficients & 15 & 1 & 155 & 67\\
Ordinary strong system & 15 & 6 & 28 & 61\\
Grouped ordinary system & 15 & 4 & 28 & 63\\
Ordinary product \(P_{69}\) & 15 & 1 & 99 & 69\\
Four-square polynomial \(Q_{56}\) & 15 & 1 & 56 & 74\\
Root-eliminated \(\Psi_{14}\) & 14 & 1 & 342 & 73\\
Root-eliminated \(\Phi_{14}\) & 14 & 1 & 293 & 75\\
Root-eliminated system & 14 & 5 & 244 & 69\\
Quadratic lift / its sum of squares & 33 & 24 / 1 & 2 / 4 & --\\
Unrestricted-integer \(\widehat P_{67}\) & \(33\;\ZZ\) & 1 & 292 & 130\\
Unrestricted-integer \(\widehat P_{68}\) & \(33\;\ZZ\) & 1 & 228 & 131\\
\bottomrule
\end{tabular}
\caption{Newly combined tradeoffs. Unmarked domains are strictly positive
integers. Degree for a system means the maximum residual degree.}
\label{tab:new}
\end{table}

Every binary addition, subtraction, and multiplication is charged, including
multiplication by a fixed numeral. Reading a fixed coefficient and reusing a
register are free. In system counts, comparisons are free; their residual
subtractions are charged when building a single polynomial. Fixed
coefficients have degree zero. Circuit registers are polynomial
abbreviations, not extra existential variables. Coefficient bit length is
reported separately because the operation metric does not charge it.

For a compact system, the four-equation version is the clearest main
recommendation. For one polynomial with fewer primary unknowns, use
\(\Psi_{14}\) or \(\Phi_{14}\). The 67-operation evaluator remains the
smallest operation count validated in this comparison, now with smaller
coefficient definitions. No result below 67 operations was established.
Classical nine-unknown existence results remain relevant: the repository's
\src{Jones1982.nine_unknowns} proves an existence statement over the naturals,
without a degree in its theorem statement~\cite{nine,jones}. Thus fourteen
is a local improvement over these delivered explicit circuits, not an
absolute variable minimum.

\Needspace{7\baselineskip}
\subsection{What was already present}

The source snapshot is pinned to commit
\begin{center}\small\texttt{3b201cdea32d2c816f5f3d20734051f25fd3ef14}.\end{center}
The eight packages had themselves been developed against the earlier pin
\begin{center}\small\texttt{cdd582bec13f908df3711fa9e05d1ddac368a5db}.\end{center}
The intake catalogue
at the newer pin carefully distinguishes mathematical proofs from finite
audits and proof-assistant builds~\cite{catalogue}. We preserve those
distinctions.

\begin{table}[ht]
\centering\small
\begin{tabular}{@{}lp{75mm}@{}}
\toprule
Package & Relevant contribution already present\\
\midrule
01 explicit82 & Fixed ordinary-input specialization; 18-unknown, 82-op polynomial.\\
02 system63 & Ordinary strong Pell relation; 16 unknowns, seven equations, degree 28.\\
03 system64 & Related 16-unknown system; root elimination yielding 15 unknowns.\\
04 poly73 & Affine marker guard and a two-sided machine normalization.\\
05 fixed15--16 & Explicit quadratic-norm root elimination; 15-unknown 80/82-op choices.\\
06 blank67 & Raw K215, no ordinary input; 15 unknowns, six equations, 61 ops; polynomial 67.\\
07 fixed70 & Signed factors and asymmetric scale choices; another low-degree system.\\
08 start-only67 & Same 67-op arithmetic skeleton as 06; a 34-unknown quartic lift.\\
\bottomrule
\end{tabular}
\caption{The existing ideas combined here. Detailed original metrics and
credits remain in the pinned catalogue.}
\end{table}

Packages 06 and 08 already give the same 67-operation formula up to
renaming, but use different coefficient integers. Packages 03 and 05 also
share a 15-unknown root-elimination construction up to an affine slack
change. We do not count either overlap as a new result. The changes here
are the \emph{combinations}: ordinary strong with inputless arithmetic;
unit-sign grouping; inputless root elimination; a wider signed-coordinate
domain; and a smaller window alphabet. Coefficients from different
compilers must never be interchanged without proving the corresponding
compiler theorem.

\section{A common arithmetic interface}\label{sec:interface}

\subsection{Fixed ports and primary unknowns}

Write the four fixed integer ports as
\[
 \beta=B-1,\qquad K=D_C+BD_R,\qquad M_C,\qquad
 M_F=M_F^{\mathrm{nat}}+B-1.
\]
In particular, \(M_F\) is the \emph{shifted} source port. The recipe gives
\(B=2^d\), \(d\ge4\), and
\begin{gather}
 0<M_C,M_F^{\mathrm{nat}}<B-1,\qquad
 M_C\equiv2\pmod4,\quad M_F^{\mathrm{nat}}\equiv4\pmod8,\label{eq:masks}\\
 \pc(M_C)+\pc(M_F^{\mathrm{nat}})=d.\nonumber
\end{gather}
These numerical conditions are useful consequences of the full recipe;
they do not by themselves characterize valid compiler coefficients.

Initially let
\[
 J,F,\alpha,z,f,h,i,T,s,w,\tau,\eta,\zeta,y,\gamma\in\Pos.
\]
Here \(z\) is the transport quotient and \(T\) the auxiliary quotient.
Define the following polynomial abbreviations:
\begin{align}
 q&=\beta J+1,& C&=q-F-\alpha,\nonumber\\
 X&=wq,&Y&=sq^3,&E&=XY,\nonumber\\
 k&=\eta+\zeta,&c&=kY+\eta,&a&=E+Y,\nonumber\\
 H&=4a+3,&\Delta&=a^2+H=(a+1)(a+3),\label{eq:core}\\
 D&=X+ac+\gamma H,&L&=E(kY),\nonumber\\
 G&=(q-1)(q-F)+\alpha=q(q-F)-C,\nonumber\\
 R&=G(q^2-1)+(M_C+qM_F)J,\nonumber\\
 V&=c(Tf-1)-Rf^2.\nonumber
\end{align}
Fixed exponents such as \(q^3\) are ordinary polynomial operations.
There are no variable exponents in any final circuit.

Four factors are common to all constructions:
\begin{align}
 N_0&=\tau^2-L(L+k),&N_1&=D^2-\Delta c^2,\nonumber\\
 N_k&=k-hE-R,&N_t&=(K+w)C+q-F-z(q-1).\label{eq:fourfactors}
\end{align}
For the inherited normalized strong block set
\begin{equation}
 S=\Delta ic^2,\quad Q=S^2,\quad
 N_a^{\rm norm}=Q(V^2-y^2)+y^2,\quad
 N_s=f^2-\Delta i^2c^4.\label{eq:normalized}
\end{equation}
The existing 67-operation polynomial is
\begin{equation}
 P_{67}=N_0N_1N_a^{\rm norm}N_kN_t(\Delta f^2-Q)-\Delta.
 \label{eq:p67}
\end{equation}
The identity
\[
 P_{67}=\Delta(N_0N_1N_a^{\rm norm}N_kN_tN_s-1)
\]
holds in the polynomial ring. On the positive domain \(\Delta>0\), so
every normalized factor at a zero is an integer unit. The first,
main, auxiliary, and strong norms cannot be \(-1\); the index and
transport signs are subsequently forced positive. The proof order matters:
transport initially gives only \(C\ge0\). The boundary
\(C=0,F=2,z=1,N_t=-1\) is not excluded until index recovery.

The inherited coefficient-reduced companion \(P_{68}\) replaces only the
auxiliary coefficient \(Q\) by \(\Delta(f^2-1)\), retaining \(Q\) in
the scaled strong factor. It has degree 123. Every zero forces \(N_s=1\),
so these coefficients agree at zeros. This is a zero-set argument, not a
polynomial identity obtained by substituting an equation off its zero set.

\subsection{A kernel fact used below}

For an integer \(A\ge2\), write
\[
 \chi_A(n)+\psi_A(n)\sqrt{A^2-1}
   =(A+\sqrt{A^2-1})^n.
\]
The repository's normalized and ordinary strong-rank lemmas, the
half-binomial Pell ratio, and the native mask/convolution compiler are the
inherited mathematical inputs~\cite{relaxed,normalized,masks,src06}.
The new arguments below explain exactly where each is used. In particular,
the signed-quotient proof is reproduced far enough to justify the
root-eliminated polynomial's positive domain; no sign of a reconstructed
root is assumed without proof.

\input{sections/ordinary_section.tex}

\section{Removing the strong root: fourteen positive unknowns}\label{sec:root14}

Delete \(f\) from the fifteen-variable list, leaving
\[
 J,F,\alpha,z,h,i,T,s,w,\tau,\eta,\zeta,y,\gamma\in\Pos.
\]
All definitions in~\eqref{eq:core} except \(V\) remain unchanged, as do
the four factors~\eqref{eq:fourfactors}. Introduce polynomial abbreviations
\begin{align}
 u&=ic^2,& S&=\Delta u,&Q&=S^2,&\mathcal D&=1+uS,\nonumber\\
 v&=cT,& b_*&=c+R\mathcal D,& P_4&=N_0N_1N_kN_t,\label{eq:rootdata}\\
 \mathcal A&=Q(\mathcal Dv^2+b_*^2-y^2)+y^2-1,\nonumber\\
 \mathcal B&=2Qvb_*,&\mathcal K&=\mathcal A^2-\mathcal D\mathcal B^2.\nonumber
\end{align}
The letter \(b_*\) here is a root-elimination expression, unrelated to the
fixed radix width \(b\) in Section~\ref{sec:coefficients}.

\begin{theorem}[Two explicit fourteen-unknown equations]\label{thm:root14}
The following ordinary integer-coefficient polynomials have a positive
zero exactly when K215 halts on blank tape:
\begin{align}
 \boxed{\Phi_{14}=P_4(1+\mathcal K)-1=0},\label{eq:phi14}\\
 \boxed{\Psi_{14}=[P_4(\mathcal A+1)-1]^2
                   -\mathcal D(P_4\mathcal B)^2=0}.\label{eq:psi14}
\end{align}
The supplied circuits use respectively \(75=44M+31A\) and
\(73=44M+29A\) operations, of exact degrees 293 and 342.
\end{theorem}

This applies the quadratic norm of packages 03/05 to the smaller inputless
base of 06/08~\cite{src05,src06,src08}. The formula for \(\mathcal B\) specifies an integer
polynomial; the circuit shares products and need not explicitly compute
that abbreviation. For example,
\[
 \mathcal D\mathcal B^2=4(Q\mathcal Dv^2)(Qb_*^2).
\]

\begin{proof}[Soundness of \(\Phi_{14}\)]
At a positive zero, \(P_4\) and \(1+\mathcal K\) are integer units.
Since \(\mathcal B\) is even,
\[
 1+\mathcal K\equiv1+\mathcal A^2\pmod4,
\]
so this factor cannot be \(-1\). Consequently \(P_4=1\) and
\(\mathcal A^2=\mathcal D\mathcal B^2\).

The denominator \(\mathcal B\) is nonzero before any index decoding.
Indeed \(Q,c,T>0\), and
\(\mathcal D=1+\Delta i^2c^4>c>0\). If the integer expression
\(c+R\mathcal D\) vanished, \(\mathcal D\) would divide \(c\), which
is impossible. Thus
\[
 f=\mathcal A/\mathcal B\in\QQ,\qquad f^2=\mathcal D\in\Pos.
\]
A rational number with integral square is integral, by reducing its
numerator and denominator to coprime terms. Thus \(f\) is a nonzero
integer. With \(V=vf-b_*\), direct expansion gives
\begin{align*}
 f^2-\Delta i^2c^4&=1,\\
 QV^2-(Q-1)y^2&=\mathcal A+1-\mathcal Bf=1.
\end{align*}
These equalities and \(P_4=1\) restore a zero of~\eqref{eq:p67},
provisionally allowing \(f\) to be signed. The signed-domain theorem of
Section~\ref{sec:signed} proves \(fT>0\); because the retained \(T\) is
positive, the restored \(f\) is positive. The inputless compiler theorem
therefore supplies an actual blank halting computation.
\end{proof}

\begin{proof}[Soundness of \(\Psi_{14}\)]
Expansion of~\eqref{eq:psi14} shows that its constant term 1 is divisible
by \(P_4\). Hence \(P_4=\pm1\), and
\[
 f=\frac{P_4(\mathcal A+1)-1}{P_4\mathcal B}
   =\frac{\mathcal A+1-P_4}{\mathcal B}
\]
is a defined rational number whose square is \(\mathcal D\). It is again
a nonzero integer. The restored auxiliary norm is
\[
 QV^2-(Q-1)y^2=\mathcal A+1-\mathcal Bf=P_4.
\]
Since \(Q\) is a square, this norm is congruent modulo four either to
\(y^2\) or to \(V^2\); it cannot be \(-1\). Therefore
\(P_4=1\), and the same restoration and signed-domain proof apply.
\end{proof}

Conversely, at any positive zero of \(P_{67}\), all six normalized
factors are one. Thus \(f^2=\mathcal D\),
\(\mathcal A=\mathcal Bf\), \(\mathcal K=0\), and \(P_4=1\).
Deleting \(f\) gives a zero of both new polynomials. The inverse is
unique because \(\mathcal B\ne0\). In particular, the two new equations
have identical positive zero sets in their fourteen retained coordinates.
The useful all-ring identity
\begin{equation}
 \Psi_{14}=P_4\Phi_{14}+(P_4-1)(2P_4\mathcal A-1)
 \label{eq:rootidentity}
\end{equation}
is included in the algebra audit, but its identity alone would not prove
the zero-set equivalence without the integer-unit argument.

\subsection{A five-equation form}

The equivalent system
\[
 N_0=N_1=N_k=N_t=1,\qquad \mathcal K=0
\]
has 14 positive unknowns, five equations, and a 69-operation shared
evaluator. Its residual degrees are \(22,18,7,2,244\). It is useful when
variable or equation count is the objective; the four-equation ordinary
system has a substantially lower degree.

\subsection{Exact degree and coefficient certificates}

Put \(q_0=\beta J\), \(k_0=\eta+\zeta\), and
\(t_0=w(q_0-F-\alpha)-zq_0\). Highest homogeneous parts give
\begin{gather*}
 (P_4)_{49}=8\gamma h k_0^3w^5s^8q_0^{29}t_0,\\
 (\mathcal D)_{34}=i^2k_0^4w^2s^6q_0^{20},\qquad
 (b_*)_{38}=i^2k_0^4w^2s^6q_0^{23}(q_0-F),\\
 Q_{46}=i^2k_0^4w^4s^8q_0^{28},\qquad
 (\mathcal A)_{122}=i^6k_0^{12}w^8s^{20}q_0^{74}(q_0-F)^2.
\end{gather*}
Here \(\deg(\mathcal D\mathcal B^2)=214<244=\deg(\mathcal A^2)\).
Therefore
\begin{align*}
 (\Phi_{14})_{293}
  &=8\gamma h i^{12}k_0^{27}w^{21}s^{48}q_0^{177}(q_0-F)^4t_0,\\
 (\Psi_{14})_{342}
  &=64\gamma^2h^2i^{12}k_0^{30}w^{26}s^{56}
         q_0^{206}(q_0-F)^4t_0^2.
\end{align*}
Both are nonzero for every \(\beta>0\). For example the coefficient of
\(J^{182}\gamma h i^{12}\eta^{27}w^{21}s^{48}z\) in \(\Phi_{14}\)
is \(-8\beta^{182}\), and that of
\(J^{212}\gamma^2h^2i^{12}\eta^{30}w^{26}s^{56}z^2\) in \(\Psi_{14}\)
is \(64\beta^{212}\). These are uniform exact-degree certificates.

\section{A larger signed domain and thirty-three integer unknowns}\label{sec:signed}

The root-elimination proof needs a domain theorem, not a choice of sign
for a square root. The following statement also improves the independent
unrestricted-integer construction.

\begin{theorem}[Six coordinates may be signed]\label{thm:signed}
For either \(P_{67}\) or \(P_{68}\), allow
\[
 \tau,y,w,h,f,T\in\ZZ
\]
while keeping \(J,F,\alpha,z,i,s,\eta,\zeta,\gamma\) strictly positive.
At every zero, \(w,h>0\), \(\tau yf\ne0\), and \(fT>0\).
Replacing \((\tau,y,f,T)\) by
\((|\tau|,|y|,|f|,\operatorname{sgn}(f)T)\) restores a positive zero.
\end{theorem}

We adapt the signed-quotient proof~\cite{signedabs,signedsound} to the
inputless arithmetic and add an elementary argument for signed \(h\).
The complete proof is also provided as \src{proofs/signed_domain.md}.

\subsection{Automatic positivity before rank recovery}

Because \(q\ge B\ge16\), \(Y=sq^3\ge q^3\), and
\(a=Y(wq+1)\), integer \(w\) gives either \(a\ge Y\) or
\(a\le-Y(q-1)\). Thus \(\Delta=(a+1)(a+3)>0\), even before
\(w\) is known positive. The normalized product has integer-unit factors.

If \(w\le-1\), then \(X,a,H<0\) and
\(D=X+ac+\gamma H<ac<0\). Hence
\(N_1> -Hc^2>1\), contradicting its being a unit. If \(w=0\),
\[
 N_1=H(2ac\gamma+\gamma^2H-c^2),\qquad H=4Y+3>1,
\]
again impossible for a unit. Consequently \(w>0\).

The negative first norm is excluded by the consecutive-product Pell
descent. For the main and strong norms, \(\Delta\equiv0,3\pmod4\)
excludes \(-1\); the auxiliary norm is congruent to a square modulo four.
All four norms are therefore \(+1\). This excludes \(\tau,y,f=0\).
Apply the absolute-value and simultaneous \(f,T\) replacement; only
\(f^2\) and \(Tf\) occur, so all polynomial values are unchanged.
We can now take \(f,\tau,y>0\), keeping \(T\) signed.
For \(P_{68}\), the strong unit first gives \(N_s=1\), after which its
auxiliary coefficient equals \(Q\); the same argument applies.

Write \(N_k=\epsilon\), \(N_t=\nu\), where
\(\epsilon\nu=1\). The transport expression is
\[
 N_t=(K+w)C+1-F-(z-1)(q-1).
\]
Its being a unit forces \(C\ge0\). The actual mask bounds then give
\begin{equation}
 3q+3<R<q^4-q^3+q^2-1<q^4-2,
 \qquad E\ge q^4>R+2.\label{eq:weakbounds}
\end{equation}
The index factor gives \(k=hE+R+\epsilon\). A negative \(h\) makes
\(k<0\), impossible. For \(h=0\), one would have \(2<k<q^4\).
But the first Pell norm classifies
\[
 k=2\psi_{P}(n),\qquad P=2XY^2+1,\quad n\ge1.
\]
The value for \(n=1\) is 2, while \(n\ge2\) gives
\(k\ge4P>q^4\). Hence \(h>0\).

\subsection{Signed rank and the index signs}

Put \(A=a+2\), so \(\Delta=A^2-1\). The main norm gives
\(D=\chi_A(p)\), \(c=\psi_A(p)\). Since \(P>A\) and \(c>k\),
\(p>n\). As \(P\equiv1\pmod E\),
\[
 2n=R+\epsilon+v_0E,\qquad v_0\ge0.
\]
Thus \(p>n\ge(R-1)/2\ge24\), \(c>2p\), and the positive index
quotient implies \(c>2R\).

The normalized strong norm gives
\(f=\chi_A(m)\), \(\psi_A(m)=ic^2\).
Strong divisibility of the Pell sequence implies \(p\mid m\). Writing
\(m=pj\), its addition formula yields
\[
 \psi_A(pj)/c\equiv jD^{j-1}\pmod c.
\]
The left side is divisible by \(c\), and \(\gcd(D,c)=1\), so
\(c\mid j\). In particular \(c\mid m\) and \(m>2p\).

The auxiliary norm classifies an odd index \(\ell=2v_1+1\) and a sign
\(e\in\{\pm1\}\) such that
\[
 V=e\,C_{v_1}(S^2),\quad
 C_{v_1}(0)=(-1)^{v_1}\ell,\quad
 C_{v_1}(-\Delta)=(-1)^{v_1}\psi_A(\ell).
\]
Here \(C_j(s^2)=\chi_s(2j+1)/s\) is an integer polynomial in \(s^2\).
The literal quotient expression, for either sign of \(T\), gives
\[
 V\equiv-c\pmod f,\qquad V\equiv-R\pmod c.
\]
As \(S^2\equiv-\Delta\pmod f\), squaring the first congruence and
using the Pell doubling identity gives
\(\chi_A(2\ell)\equiv\chi_A(2p)\pmod{\chi_A(m)}\).
The plus-sign step-down gives
\(2\ell\equiv\pm2p\pmod{4m}\). One proves this by reducing the
index modulo \(2m\), reflecting it into \([0,m]\), and using
\(2\chi_A(m-1)<\chi_A(m)\) to exclude a negative congruence between
positive representatives. The endpoint \(m\) has residue zero and is
also excluded. The second congruence, and \(c\mid S,m\), now give
\(R\equiv\pm p\pmod c\). Since \(R,p\in(0,c/2)\), \(p=R\).

The no-wrap bound~\eqref{eq:weakbounds} forces \(v_0=0\). If
\(\epsilon=-1\), then \(p=2n+1\), whereas Pell duplication gives
\[
 \psi_A(2n)=2A\psi_{2A^2-1}(n)\ge Ak,
\]
because \(2A^2-1>P\). This contradicts
\(c=\psi_A(2n+1)>Ak>k(Y+1)\). Thus both index and transport units are
\(+1\); in particular \(C>0\), and \(R\) is odd.

\subsection{Power recovery, masks, and the final sign}

Write \(r=(R-1)/2\ge24\) and
\(\xi=(X+1)^{2r}/X^r\). The elementary Pell estimates give
\(c/(k/2)>\xi\), hence \(Y>X^r/3\) and \(a>2^R\).
The projection recurrence gives
\(\chi_A(R)-a\psi_A(R)\equiv2^R\pmod H\).
The main-root definition therefore implies \(X\equiv2^R\pmod H\).
Both positive quantities are below \(a<H\), so \(X=2^R\).
Consequently \(q=2^t=B^N\), and \(q^3\mid X\).

The upper ratio estimate uses only odd \(R\), not a positive \(V\):
\[
 0<c/k-\xi/2<16r/(X+1)<1/2.
\]
Writing
\[
 \xi=M+\theta,\quad
 M=\binom{2r}{r}+\sum_{j=1}^{r}\binom{2r}{r+j}X^j,
 \quad0<\theta<1/4,
\]
one finds that \(M\) is even and \(Y=M/2\). The exact binomial
valuation gives \(v_2(Y)=\pc(r)-1\), so \(q^3\mid Y\) implies
\(\pc(R)\ge3t+2\).

Define the finite binary words
\[
 S_w=C-1+qF,\qquad
 T_w=(M_CJ+1)+q(M_F^{\mathrm{nat}}J-1).
\]
The packing identity is
\[
 R=(q^2-S_w)(q^2-1)+T_w,
 \quad0<S_w<q^2,\quad\pc(T_w)=t+2.
\]
The inverse-population lemma gives \(\pc(R)\le3t+2\), with equality
exactly when \(S_w\mathbin{\&}T_w=0\). Thus
\((C-1)\mathbin{\&}(M_CJ+1)=0\).
Because \(M_CJ+1\equiv3\pmod4\), this forces \(C\equiv1\pmod4\),
and the literal packing identity gives \(R\equiv3\pmod4\).

The stronger step-down above has
\(\ell=\epsilon_0R+2mj_0\), \(\epsilon_0\in\{\pm1\}\).
Substitution into \emph{both} unsquared auxiliary congruences gives
\begin{align*}
 e(-1)^{r+mj_0}R&\equiv-R\pmod c,\\
 e(-1)^{r+(m+1)j_0}c&\equiv-c\pmod f.
\end{align*}
As \(c>2R\) and \(f>2c\), both signs must equal \(-1\). Their
quotient forces \(j_0\) even. Since \(r\) is odd, \(e=1\), so \(V>0\).
Finally \(cTf=V+c+Rf^2>0\) proves \(T>0\) in the transformed tuple,
or \(fT>0\) before transformation. This proves Theorem~\ref{thm:signed}
without applying a positive-domain compiler to a signed tuple.

\subsection{Explicit conversion to unrestricted integers}

Every positive integer is represented by
\[
 \operatorname{Pos}(u,v,t)=1+u^2+v^2+t(t+1),\qquad u,v,t\in\ZZ.
\]
This is always positive. Surjectivity follows from the three-square
theorem applied to \(4n-3\): its residue modulo four forces exactly one
odd square, giving the displayed formula. Substitute a disjoint copy for
each of the nine still-positive coordinates and keep the other six signed.
The result has \(9\cdot3+6=33\) unrestricted integer unknowns. Each
substitution costs seven operations, so the counts are 130 and 131.

The exact degrees are smaller than naive degree doubling. Giving weight
two to the nine substituted coordinates and weight one to the six signed
ones, the \(P_{67}\) factor degrees are
\(42,34,112,12,4,88\), totaling 292. For \(P_{68}\), the auxiliary
degree is 48, giving total 228. The nonzero top forms are respectively
\begin{gather*}
 8q_0^{92}h\gamma i^4k_0^{13}w^{13}s^{26}T^2f^2z,\\
 8q_0^{72}h\gamma i^2k_0^9w^{11}s^{20}T^2f^4z.
\end{gather*}
Substitution of the disjoint positive quadratic forms preserves
nonvanishing: setting the second and third new coordinates to zero gives
the injective monomial substitution \(x\mapsto u_x^2\).

\input{sections/projected_cross_section.tex}

\section{Why ZFC proves the equivalence}\label{sec:logic}

The literal machine is the 215-state binary K215 table whose SHA-256 is
\begin{center}\scriptsize\texttt{fdbcef0825886035c3091c8537a69535bc96f9e265f0ae45863543d40ae70090}.\end{center}
The inspected endpoint has the statements
\begin{quote}\small\ttfamily
KEndpoint215.blank\_halts\_iff\_zfc\_refutation\\
KEndpoint215.nonhalting\_iff\_zfc\_consistency.
\end{quote}
They use the syntactic predicate \texttt{BProv ZFCax\_s [] Form.fBot}.
The machine-specific proof invokes the common K endpoint and the
constructible-choice equivalence between finite ZF and ZFC refutations;
it is not an assertion that consistency is true~\cite{k215}.

The resulting ordinary mathematical chain is
\begin{align*}
 \exists\text{ Diophantine solution}
 &\ \Longleftrightarrow\ \exists\text{ finite accepting local tableau}\\
 &\ \Longleftrightarrow\ \text{K215 halts on blank tape}\\
 &\ \Longleftrightarrow\ \exists p\,\operatorname{Proof}_{\mathrm{ZFC}}(p,\bot).
\end{align*}
Negating this equivalence yields the theorem in the requested direction.
The finite table, coefficient algorithms, polynomial manipulations,
integer divisibility arguments, Pell recurrences, and finite-computation
calibration can all be represented in ZFC. The mathematical argument does
not require ZFC to prove that Lean or Rocq is sound. Nor does it infer a
ZFC-internal proof merely from an external proof-assistant theorem:
the intended internal proof is the ordinary syntactic calibration together
with the explicitly described arithmetical reduction.

\paragraph{Formal status.}
This packet provides an ordinary proof using the pinned compiler/Pell
development and the new lemmas proved here. It does not contain a new
machine-checked term for the complete Diophantine theorem, or an encoded
derivation in ZFC's own proof syntax. Lean and Rocq endpoint source and
historical verification records were inspected; their full builds were
not rerun. The native compiler itself is a mathematical research
development, not a completed Lean compiler certificate. These are material
verification boundaries, shared with the eight predecessor packages.

\paragraph{Incompleteness.}
If ZFC is consistent, the displayed equation really has no solution, but
ZFC cannot prove that unsolvability statement without proving its own
consistency. Mere consistency does not justify a claim that ZFC also
cannot prove the opposite existential assertion; a corresponding
soundness assumption is needed for that stronger two-sided independence
claim. No consistency assumption is used in constructing the equivalence.

\section{Reproducibility and what the checks establish}\label{sec:checks}

The archive contains the literal K215 table, the four-port coefficient
specification, every arithmetic instruction, the quadratic residuals and
expanded quartic, proof notes, inherited evidence, and provenance hashes.
Run
\begin{quote}\ttfamily python3 code/verify.py\end{quote}
to reproduce the bounded checks with Python's standard library. Rebuild
this report with \texttt{make pdf}. The report's TeX and all section files
are included. No network access is needed by verification.
The default run checks and reuses the delivered exact local-predicate
census; use \texttt{python3 code/verify.py -{}-full} to rerun all thirteen
center classes and their left/right pairs.

The checks cover the following independently auditable claims:
\begin{enumerate}
\item K215's exact table hash and the finite local-predicate census;
projected alphabet counts by exact fiber sizes, rather than enumeration
of all projected tiles.
\item Actual-dimension interval support incidences, clause-band separation,
anchor differences, mask population formulas, and the fixed modulo-five
unit; smaller sparse examples exercise the same coefficient construction.
\item Acyclic source binding, every live gate and port, exact operation
counts, and exact source/formula identities on signed rational or integer
diagnostic assignments.
\item Exact sparse expansion of the 24 quadratic equations and their
quartic, including the coefficient 1 on \(\tau^4\).
\item Exact polynomial evaluation on generic lines over a finite field,
corroborating degrees whose uniform upper bounds and nonzero leading
forms are proved in the text.
\end{enumerate}
None of these finite checks is a test of consistency, a proof by sampling,
or an exhibited K215 halting computation. The enormous coefficient integers
are not materialized. The output is fully specified by finite expressions
and algorithms; it is not a practical expanded decimal polynomial.

\subsection{Recommended incorporation}

Retain the four-equation system and \(Q_{56}\) as the readable low-degree
entry point, \(\Psi_{14}\) and \(\Phi_{14}\) as the fourteen-variable
alternatives, and the projected-cross recipe as an independent coefficient
improvement. Keep the 67-operation circuit as a baseline. The source
comments and proof notes distinguish the normalized strong coordinate
from the ordinary strong coordinate: the conversion is an existence
equivalence, and arbitrary ordinary solutions must not be divided by
\(\Delta\) without the relaxed-rank proof.

The next consequential formalization task is the common inputless
compiler interface plus the signed-domain theorem. That would certify
several of these tradeoffs at once. A search for fewer operations should
work from the complete four-port circuit and its domain proof; deleting a
Pell or population condition merely because it holds on a few examples
does not establish a valid reduction.

\begin{thebibliography}{99}\small
\bibitem{catalogue} ProveIt contributors. \emph{Con(ZFC) in Diophantine form},
 eight-package catalogue, pinned 22 October 2026.
 \href{https://github.com/VladimirReshetnikov/ProveIt/blob/3b201cdea32d2c816f5f3d20734051f25fd3ef14/SetTheory/Cardinals/docs/reports/hilbert-tenth-problem/con-zfc-diophantine-forms/README.md}{Repository source}.
\bibitem{src02} ProveIt contributors. \emph{A fixed Diophantine representation of Con(ZFC)},
 package 02-system63, report and ordinary strong proofs, same pin.
\bibitem{src05} ProveIt contributors. \emph{Fixed-input Diophantine forms of ZFC consistency},
 package 05-fixed15-16, Sections 4--7 and emitted circuits, same pin.
\bibitem{src06} ProveIt contributors. \emph{Explicit Diophantine equations for Con(ZFC)},
 package 06-blank67, report, fixed-blank proof, helical proof, and coefficient compiler, same pin.
\bibitem{src08} ProveIt contributors. \emph{Explicit Diophantine equations for the consistency of ZFC},
 package 08-start-only67, including \src{start_only_quartic34_proof.md}, same pin.
\bibitem{relaxed} ProveIt contributors. \emph{Pell relaxed auxiliary proof},
 \src{Computability/HilbertTenthProblem/Papers/1980/PELL_RELAXED_AUXILIARY_PROOF.md}, same pin.
\bibitem{normalized} ProveIt contributors. \emph{An auxiliary quotient gives a complete universal polynomial in 85 operations},
 \src{complete85_auxiliary_bezout_projection.md}, native-stream-queue research notes, same pin.
\bibitem{masks} ProveIt contributors. \emph{Modified native masks for the half-binomial complete75 compiler},
 \src{complete75_half_binomial_compiler.md}, same pin.
\bibitem{signedabs} ProveIt contributors. \emph{Signed auxiliary quotient absorption},
 \src{complete84_signed_quotient_absorption.md}, especially Section 2, same pin.
\bibitem{signedsound} ProveIt contributors. \emph{Signed auxiliary quotient: source-specific recovery of the positive domain},
 \src{complete84_signed_quotient_soundness_scout.md}, and its independent review, same pin.
\bibitem{k215} ProveIt contributors. K215 endpoint and verification records.
 \href{https://github.com/VladimirReshetnikov/ProveIt/blob/3b201cdea32d2c816f5f3d20734051f25fd3ef14/Computability/BusyBeaver/Lean/BusyBeaver/ZFC/KEndpoint215.lean}{KEndpoint215.lean};
 \href{https://github.com/VladimirReshetnikov/ProveIt/tree/3b201cdea32d2c816f5f3d20734051f25fd3ef14/Computability/BusyBeaver/Research/ZFC/k225-formal/verification/k215}{verification directory}.
\bibitem{nine} ProveIt contributors. \src{Jones1982.nine_unknowns},
 \src{Computability/HilbertTenthProblem/Lean/Diophantine/Paper1982/NineMain.lean}, same pin.
\bibitem{jones} J. P. Jones. Universal Diophantine equation.
 \emph{Journal of Symbolic Logic} 47(3), 549--571 (1982).
 \href{https://doi.org/10.2307/2273588}{doi:10.2307/2273588}.
\end{thebibliography}
\end{document}

